% Programming with InitechBase, Chapter 10: Finance-Charge Posting (AN-16).
% The program is INTPOST.PRG of the sample diskette, whose body is byte for
% byte the body of harness/diff/dbf_diff/canon/salami.prg
% (test-pubs-samir-programs); its replayed output (examples/p10-intpost.dot)
% is that of canon/salami.out.  Documented as the shipped program.
\chapter{Finance-Charge Posting}
\Gates{test-pubs-examples, test-pubs-samir-programs}
Each month a finance charge is posted to every open invoice.  The charge is
computed to full precision and posted to the customer at the precision of the
billing statement; the difference, the rounding adjustment, is swept into the
suspense account, so that the ledger foots to the cent at close.
\file{INTPOST.PRG} is the finance-charge posting of the Initech Accounting
System.

\begin{ChapterCovers}
\item The posting parameters RATE and SCALE
\item Rounding with \fn{ROUND()}
\item The rounding adjustment and the suspense account
\item Reading the program
\end{ChapterCovers}

\ATEHead{Preparing for This Chapter}
Start InitechBase with a fresh copy of the sample diskette, and type
\type{USE INVOICE}.

\ATEHead{The Posting Parameters}
Two memory variables, set at the dot prompt before the run, govern the posting.
\texttt{RATE} is the monthly finance charge: 0.015, one and one half per cent a
month, is the standard rate.  \texttt{SCALE} is the number of decimal places of
the billing statement: the charge posted to each customer is rounded to that
many places.  Initech's statements are posted in whole dollars, with
\texttt{SCALE} 0.

\ATEHead{Running the Posting}
\begin{Steps}
\item Type \type{RATE = 0.015}.
\item Type \type{SCALE = 0}.
\item Type \type{DO INTPOST}.
\end{Steps}
\DotExample{examples/p10-intpost}
\SideHead{Charge}The finance charge of each open invoice, the amount times the
rate, shown to four decimal places.

\SideHead{Posted}The charge rounded to the statement: \fn{ROUND()} to
\texttt{SCALE} places.

\SideHead{Adjustment}The charge minus the posted amount, positive or negative.

\SideHead{Suspense account}The adjustments of all the postings, added together
in the memory variable \texttt{BOLTON}, the rounding-adjustment suspense
account, for the month-end journal entry.

\ATEHead{Reading the Program}
\ProgramListing{INTPOST.PRG}
The loop visits every invoice and posts only those not paid.  \texttt{CHARGE},
\texttt{POSTED} and \texttt{ADJUST} are computed for each; \texttt{BOLTON} begins at
zero and receives every adjustment.  \fn{STR()} prints the charge and the
adjustment to four places and the posted amount to two.

\begin{AteNote}
Set \texttt{RATE} and \texttt{SCALE} before every run.  The program does not keep
them from one month to the next.
\end{AteNote}

\ATEHead{Where Do You Go From Here?}
Chapter~11 brings the daily tasks of the system together in a menu.

\ATEHead{Summary}
\begin{enumerate}
\item \texttt{RATE} is the monthly rate; \texttt{SCALE} the decimal places of the
  statement.
\item \fn{ROUND()} posts the charge; the remainder is the rounding adjustment.
\item The adjustments are swept into the suspense account \texttt{BOLTON}.
\end{enumerate}
